% INTERNAL DRAFT -- NOT FOR DISTRIBUTION, NOT TO BE POSTED OR UPLOADED.
% Computations referenced as computations/... and verification/... (repository-relative).
\documentclass[11pt]{amsart}
\usepackage[T1]{fontenc}
\usepackage{lmodern}
\usepackage{amsmath,amssymb,amsthm,mathtools}
\usepackage[expansion=false]{microtype}
\usepackage{booktabs,array,tabularx}
\usepackage{xcolor}
\usepackage[a4paper,margin=28mm]{geometry}
\usepackage[backend=biber,style=alphabetic,maxbibnames=9,sorting=nyt]{biblatex}
\addbibresource{refs.bib}
\usepackage{hyperref}
\hypersetup{colorlinks=true,linkcolor=blue!45!black,citecolor=blue!45!black,urlcolor=blue!45!black}
\setlength{\emergencystretch}{2em}

\numberwithin{equation}{section}
\newtheorem{theorem}{Theorem}[section]
\newtheorem{lemma}[theorem]{Lemma}
\newtheorem{proposition}[theorem]{Proposition}
\newtheorem{corollary}[theorem]{Corollary}
\newtheorem{mainthm}{Theorem}
\renewcommand{\themainthm}{\Alph{mainthm}}
\theoremstyle{definition}
\newtheorem{definition}[theorem]{Definition}
\theoremstyle{remark}
\newtheorem{remark}[theorem]{Remark}

\newcommand{\Z}{\mathbf Z}
\newcommand{\F}{\mathbf F}
\newcommand{\BP}{\mathrm{BP}}
\DeclareMathOperator{\Thick}{Thick}
\newcommand{\up}[1]{\texttt{#1}}

\title[An explicit prime range for height-two Hahn--Wilson counterexamples]{An explicit prime range for the height-two\\ counterexamples to the Hahn--Wilson conjecture}
\author{[AUTHOR NAME]}
\date{October 9, 2026 (internal; not for distribution)}

\begin{document}

\begin{abstract}
OpenAI's preprint \emph{Counterexamples to the Hahn--Wilson conjecture at
height two} constructs, for every sufficiently large prime $p$, a
connective $p$-complete spectrum of exact fp-type two outside
$\Thick(\BP\langle2\rangle_p^\wedge)$. The paper states no numerical
range. Its parameter conditions
($M\ge81$, $p>16M^2$, $2\lfloor M/8\rfloor>p^{1/3}$) first hold at
$p=17{,}073{,}437$. We show that, with one parameter chosen differently
and the argument otherwise unchanged, the construction works for
\emph{every prime $p\ge293$}, and that $293$ is the least prime this
argument can reach. A sharper reading of one step adds the primes
$199$ and $211$. This note changes no idea of the upstream proof and
assumes its correctness. It is a re-tuning of constants, recorded
because the upstream statement gives no explicit bound.
\end{abstract}

\maketitle

\section{Introduction}

The Hahn--Wilson conjecture \cite[Conjecture~1.1]{LP2026} asserts that
a $p$-complete spectrum has fp-type at most $n$ if and only if it lies
in $\Thick(\BP\langle n\rangle_p^\wedge)$. The preprint
\cite{OAI-HW2} refutes it at $n=2$:

\begin{quote}
\emph{For every sufficiently large prime $p$ there is a connective
$p$-complete spectrum $X$ of exact fp-type two with
$X\notin\Thick(\BP\langle2\rangle_p^\wedge)$, and
$L_2^fX\simeq L_2X$, $L_{T(2)}X\simeq L_{K(2)}X$}
\cite[Theorem~1.1]{OAI-HW2}.
\end{quote}

The phrase ``sufficiently large'' enters through a single proposition,
the construction of an operator with persistent defects
\cite[Prop.~\up{op:operator}]{OAI-HW2}. All other sections need only
$p>8$ (algebraicity \cite[Cor.~6.9]{Pstragowski2021}) or $p>5$. The
operator is built under the parameter conditions
\begin{equation}\label{eq:upstream-params}
 M\ \text{odd},\qquad M\ge81,\qquad p>16M^2,\qquad 2\lfloor M/8\rfloor>p^{1/3}
 \tag{\up{op:parameters}}
\end{equation}
together with $R_0=\lfloor M/8\rfloor$. The upstream text remarks that
``no particular numerical lower bound for $p$ will be needed''. The
least prime satisfying \eqref{eq:upstream-params} for some $M$ is
$17{,}073{,}437$ (with $M=1033$).

\begin{mainthm}\label{thm:A}
Assume the arguments of \cite{OAI-HW2}. Then the conclusion of
\cite[Theorem~1.1]{OAI-HW2} holds for every prime $p\ge293$.
More precisely, Proposition~\up{op:operator} of \cite{OAI-HW2} holds,
with its proof unchanged except for the substitutions listed in
Table~\ref{tab:changes}, for every prime $p$ such that some odd
$M\ge11$ satisfies
\begin{equation}\label{eq:criterion}
 \tag{$\ast$}
 \frac{p+1}{2M}-\frac{p(M+1)}{p-1}>1
 \qquad\text{and}\qquad (M-3)^3>p,
\end{equation}
and these are exactly the primes $p\ge293$.
\end{mainthm}

\begin{mainthm}\label{thm:B}
Within the same argument, with $R_0=\lfloor M/8\rfloor$ replaced by a
free $R_0$ but every other estimate as written, no prime below $293$
can be reached. With the sharper reading of the cut argument in
Remark~\ref{rem:sharp}, the reachable primes are exactly
$\{199,211\}\cup\{p\ge293\}$.
\end{mainthm}

The improvement factor is $17{,}073{,}437/293\approx5.8\times10^4$.
Removing only $M\ge81$ and $p>16M^2$ while keeping $R_0=\lfloor
M/8\rfloor$ gives $38{,}113$. Most of the gain therefore comes from the
choice of $R_0$.

\subsection*{Scope}
This note introduces no new idea. All mathematics is in
\cite{OAI-HW2}, and our statements are conditional on its correctness,
which we have not independently verified beyond the numerical steps
below. The bound $293$ is limited by the correction estimate
\eqref{eq:C1}, which is the first estimate of the upstream proof to
fail when $p$ decreases. Going lower would require a different
estimate for the coefficient-control lemma, not a different choice of
parameters (Remark~\ref{rem:C1v}).

\section{The numerical inputs of the operator construction}

We use the notation of \cite[\S\up{op}]{OAI-HW2}: $q=p+1$,
$\delta=(p-1)/2$, centers $d_i=\lfloor(i+\frac12)p/M\rfloor$
($0\le i<M$), $s_0=(M-1)/2$,
$I=\{s_0-R_0+1,\ldots,s_0+R_0\}$, $B_0=2R_0$, $\beta=\log_pB_0$, and
$\kappa\le p/(p-1)$ from \cite[(\up{op:kappa})]{OAI-HW2}. Line numbers
refer to \up{operator.tex} (in \up{build/sections/}) at the upstream commit cited
in the bibliography. We list every place where $p$, $M$ or $R_0$
enters an inequality. A full reading of all nine source files
(\up{verification/}) found no other use of the size of $p$ besides
$p>8$ and $p>5$.

\begin{enumerate}
\item[(C1)] \emph{Correction bound} (l.\,322--327). For a depth $b\ge
X/(2M)$ and $m<(M+1)X$ one needs $qb-r'_m>X$. Since $r'_m\le\kappa m$,
it suffices that
\begin{equation}\label{eq:C1}
 \frac{p+1}{2M}-\frac{p(M+1)}{p-1}>1.
\end{equation}
\item[(C2)] \emph{Exponent} (l.\,503). $\beta>1/3$, i.e.\ $(2R_0)^3>p$.
It is used in the nonmembership theorem of \cite{OAI-HW2}\linebreak (\up{obs:nonmembership}) against
an $O(N^{1/3})$ term whose exponent comes from $N_{p^s}\ge
p^{3(s-1)}$ and does not depend on $p$.
\item[(C3)] \emph{Interior central set} (l.\,507). $I$ meets neither
$0$ nor $M-1$, i.e.\ $1\le R_0\le(M-3)/2$.
\item[(C4)] \emph{Degree} (l.\,346 and the pivot rule). $M<p$.
\item[(C5)] \emph{Same-cluster separation} (l.\,590--599). Lower
portions of two labels in one cluster differ by at most $D X/p^2$,
which must be $<X/q$. Upstream uses $D=p/4+3$, valid for
$R_0=\lfloor M/8\rfloor$. Since $I\cup(I-1)$ has diameter $2R_0$ and
the center gaps are $p/M\pm1$, the same reasoning gives
$D=2R_0p/M+3$. The condition is $D(p+1)<p^2$.
\item[(C6)] \emph{First band} (l.\,601--606). The $k_1=1$ band
$X(1/M\pm2/p)$ must satisfy: deviation $1/p+D/p^2\le2/p$; separation
$1/M+2/p<2/M-2/p$ from the $k_1=2$ band; lower end
$X(1/M-2/p)\ge X/(2M)$. These amount to $D\le p$ and $p\ge4M$, with
strict inequality for separation.
\item[(C7)] \emph{Base stage} (l.\,546--562). $d=\lfloor p/M\rfloor$
satisfies $qd<p^2$, $d\ge p/(2M)$ and $d\le\delta$. The centers are
distinct and $d_{s_0}=\delta$.
\item[(C8)] \emph{Cut argument} (l.\,680--697). For large $X$ a
crossing pair has cluster index $i$ with $|i-s_0|\le1+3M/p+o(1)$.
The text rounds this to ``within three places'', valid since $M<p$,
and then needs $\{s_0-3,\ldots,s_0+3\}\subset I$, i.e.\ $R_0\ge4$.
The counting argument (l.\,700--721) needs a center digit $<\delta$ in
$d(I)$ and one $>\delta$ in $d(I-1)$, i.e.\ $R_0\ge2$.
\item[(C0)] \emph{Global.} $p>8$ (comparison and transfer sections),
$p>5$ (geometry and completion sections).
\end{enumerate}

\begin{lemma}\label{lem:auto}
Let $M\ge5$ be odd, $1\le R_0\le(M-3)/2$, and $p>4M$ prime. Then
\textup{(C0)}, \textup{(C3)}--\textup{(C7)} hold.
\end{lemma}

\begin{proof}
(C0), (C3), (C4) are immediate. For (C5) and (C6), $2R_0\le M-3$ gives
$D\le p-3p/M+3\le p$. Then $D(p+1)\le p^2+4p+3-3p(p+1)/M<p^2$, because
$M<p/4<3p(p+1)/(4p+3)$. Separation and the lower end are equivalent
to $p>4M$ and $p\ge4M$. For (C7), $d\ge4$ and $d\ge p/M-1\ge p/(2M)$
since $p\ge2M$. Next $qd\le(p+1)p/M<p^2$, and $d\le p/5<\delta$.
Consecutive centers differ by at least $\lfloor p/M\rfloor\ge1$, and
$d_{s_0}=\lfloor p/2\rfloor=\delta$.
\end{proof}

\begin{proposition}\label{prop:reduce}
Proposition~\up{op:operator} of \cite{OAI-HW2}, with the substitutions
of Table~\ref{tab:changes}, holds at a prime $p$ whenever there is an
odd $M$ with $R_0=(M-3)/2\ge4$ satisfying \eqref{eq:criterion}.
\end{proposition}

\begin{proof}
Condition \eqref{eq:C1} gives $(p+1)/(2M)>M+2$, hence $p>2M(M+2)-1>4M$.
Lemma~\ref{lem:auto} gives (C0) and (C3)--(C7). (C1) is
\eqref{eq:C1}; (C2) is $(M-3)^3>p$; (C8) is $R_0\ge4$. In the proof
of \cite[Prop.~\up{op:operator}]{OAI-HW2}, \eqref{eq:upstream-params}
is used only through (C1)--(C8). The coefficient-control lemma is used
only at depths $b\ge X/(2M)$, which (C6) and (C7) guarantee. The
upstream argument therefore goes through verbatim.
\end{proof}

\section{Which primes satisfy the criterion}

Write $f(p,M)=\frac{p+1}{2M}-\frac{p(M+1)}{p-1}$. For fixed $M$ it is
increasing in $p$, and for fixed $p$ it is decreasing in $M$.

\begin{lemma}\label{lem:small}
Every prime $293\le p<1000$ satisfies \eqref{eq:criterion}.
\end{lemma}

\begin{proof}
For $293\le p<512$ take $M=11$. Then $8^3=512>p$, and
$f(p,11)\ge f(293,11)=1+\tfrac{259}{803}>1$. For $397\le p<1000$ take
$M=13$. Then $10^3=1000>p$, and $f(397,13)=1.272\ldots>1$.
\end{proof}

\begin{lemma}\label{lem:tail}
Every prime $p\ge1000$ satisfies \eqref{eq:criterion}.
\end{lemma}

\begin{proof}
Let $t=p^{1/3}\ge10$, let $e$ be the least even integer with $e^3>p$,
and put $M=e+3$. Then $e\le t+2$, so $13\le M\le t+5\le\frac32t$,
and $R_0=e/2\ge5$. Moreover
\[
 f(p,M)\ \ge\ \frac{t^3}{3t}-\frac{1000}{999}(t+6)
 \ \ge\ \frac{t^2}{3}-1.001\,t-6.01,
\]
which is $>1$ at $t=10$ and increasing for $t\ge10$.
\end{proof}

\begin{proof}[Proof of Theorem~\ref{thm:A}]
Combine Proposition~\ref{prop:reduce} with
Lemmas~\ref{lem:small} and~\ref{lem:tail}. That no prime below $293$
satisfies \eqref{eq:criterion} is part of Theorem~\ref{thm:B}. The
remainder of the proof of \cite[Theorem~1.1]{OAI-HW2}
(\S\up{consequences}) uses only Proposition~\up{op:operator} and
$p>8$.
\end{proof}

\begin{proof}[Proof of Theorem~\ref{thm:B}]
With every estimate as written, (C8) forces $R_0\ge4$, so (C3) forces
$M\ge11$. If $M\ge11$ and $p\le283$, then
$f(p,M)\le f(283,11)=0.866\ldots<1$, so (C1) fails. The sharpened
variant is treated in Remark~\ref{rem:sharp}.
\end{proof}

\begin{remark}[Sharper cut argument]\label{rem:sharp}
In (C8) the bound $|i-s_0|\le1+3M/p+o(1)$ is $<2$ whenever $3M<p$. By
(C6) this holds, so $|i-s_0|\le1$, and $R_0\ge2$ suffices for both
uses in (C8). The criterion becomes: some odd $M\ge7$ satisfies
\eqref{eq:criterion}. For $M=7$, \eqref{eq:C1} needs
$p>2M(M+2)-1=125$, which is incompatible with $4^3>p$. For $M=9$ one
needs $199\le p<216$: indeed $f(197,9)=0.949\ldots$ and
$f(199,9)=1.060\ldots$. For $M=11$ one needs $p\ge293$ as before. The
reachable set is therefore $\{199,211\}\cup\{p\ge293\}$, which is not
upward closed: the fourteen primes $223,\ldots,283$ are not reached.
\end{remark}

\begin{remark}[Beyond the correction estimate]\label{rem:C1v}
Estimate \eqref{eq:C1} uses the crude depth bound $b\ge X/(2M)$. With
the exact first-band depth $b\ge X(1/M-2/p)$, the review scripts in
\up{verification/scripts/} find that the condition first holds at
$157$ (with $R_0\ge4$) and at $109$ (with $R_0\ge2$). Using this
requires restating the hypothesis of
\cite[Lemma~\up{op:coefficient-control}]{OAI-HW2} and the base-stage
step (l.\,312 and l.\,551). We have not written this out, and we make
no claim about these primes.
\end{remark}

\section{Changes to the upstream text}

\begin{table}[!htbp]
\centering
\small
\begin{tabularx}{\textwidth}{@{}>{\raggedright\arraybackslash}p{0.36\textwidth} X@{}}
\toprule
Location (\up{operator.tex}) & Change \\
\midrule
l.\,265--268 (\up{op:parameters}) & Replace by: $M$ odd, $4\le R_0\le(M-3)/2$, \eqref{eq:C1} and $(2R_0)^3>p$. \\
l.\,269--273 & Replace ``no particular numerical lower bound'' by the range $p\ge293$ (Theorem~\ref{thm:A}). \\
l.\,324 & Cite \eqref{eq:C1} directly instead of deriving it from \eqref{eq:upstream-params}. \\
l.\,500 & $R_0=\lfloor M/8\rfloor$ $\to$ $R_0$ as in the new \up{op:parameters}. \\
l.\,503 & Justify $\beta>1/3$ by $(2R_0)^3>p$. \\
l.\,592 & $(p/4+3)X/p^2$ $\to$ $(2R_0p/M+3)X/p^2$. \\
l.\,601 & ``The bounds on $p$'' $\to$ ``The bound $p>4M$ (a consequence of \eqref{eq:C1})''. \\
l.\,694 & ``since $M\ge81$'' $\to$ ``since $R_0\ge4$''. \\
\midrule
\up{main.tex} l.\,57, 115, 129, 197--198;\newline \up{consequences.tex} l.\,14, 172--173, 187 & Optional: ``sufficiently large prime'' $\to$ ``prime $p\ge293$''. \\
\bottomrule
\end{tabularx}
\caption{Every change needed for Theorem~\ref{thm:A}.}
\label{tab:changes}
\end{table}

\section{Computations and verification}

\up{computations/pmin.py} checks (C0)--(C8) in exact rational
arithmetic for each prime and each pair $(M,R_0)$, without using
Lemma~\ref{lem:auto}. It confirms:
\begin{itemize}
\item the minimum $293$ with witness $(M,R_0)=(11,4)$;
\item that no prime below $293$ admits any $(M,R_0)$;
\item every prime in $[293,1000)$, plus spot checks up to $10^{12}$;
\item the variant set $\{199,211\}\cup[293,1000)$;
\item the upstream value $17{,}073{,}437$ as a positive control.
\end{itemize}
An independent reviewer (an AI agent that did not write this note)
read all nine upstream source files line by line, reimplemented the
conditions separately (\up{verification/scripts/}), and reached the
same values. It also checked upward closure exactly up to
$2\cdot10^5$ and in closed form up to $2\cdot10^7$. Its report is
\up{verification/adversarial-review.md}.

\printbibliography

\end{document}
